% Compile with LuaLaTeX (Japanese labels).
\documentclass[border=3mm]{standalone}
\usepackage{luatexja}
\usepackage{tikz,amsmath,amssymb,tikz-feynhand,circuitikz}
\usetikzlibrary{arrows.meta}
\begin{document}
% Generated by TikZ Draw. Coordinates: mm; line widths and fonts: TeX pt.
% Requires: \usepackage{tikz,amsmath,amssymb,tikz-feynhand,circuitikz} and \usetikzlibrary{arrows.meta}
\begin{tikzpicture}[x=1mm,y=1mm,line cap=round,line join=round]
  \definecolor{td0}{HTML}{226B61}
  \definecolor{td1}{HTML}{263544}
  \path[use as bounding box] (0,0) rectangle (160,200);
  \begin{feynhand}
    \setlength{\feynhandlinesize}{0.8pt}
    \vertex (tdf0a) at (15,150);
    \vertex (tdf0b) at (50,150);
    \propag[photon,color=td0,mom'={[arrow style=td0] $q$}] (tdf0a) .. controls (26.6667,150) and (38.3333,150) .. (tdf0b) node[pos=0.5,auto=left,inner sep=2pt,text=td0] {$\text{光子}$};
  \end{feynhand}
  \begin{feynhand}
    \setlength{\feynhandlinesize}{0.8pt}
    \vertex (tdf1a) at (110,150);
    \vertex (tdf1b) at (145,150);
    \propag[photon,color=td0] (tdf1a) .. controls (121.6667,150) and (133.3333,150) .. (tdf1b) node[pos=0.5,auto=left,inner sep=2pt,text=td0] {$\gamma$};
  \end{feynhand}
  \begin{feynhand}
    \setlength{\feynhandlinesize}{0.8pt}
    \vertex (tdf2a) at (50,150);
    \vertex (tdf2b) at (110,150);
    \propag[fermion,color=td1,mom={[arrow style=td1] $k$}] (tdf2a) .. controls (50,192) and (110,192) .. (tdf2b) node[pos=0.5,auto=right,inner sep=2pt,text=td1] {$e^-$};
  \end{feynhand}
  \begin{feynhand}
    \setlength{\feynhandlinesize}{0.8pt}
    \vertex (tdf3a) at (110,150);
    \vertex (tdf3b) at (50,150);
    \propag[fermion,color=td1,mom={[arrow style=td1] $k+q$}] (tdf3a) .. controls (110,108) and (50,108) .. (tdf3b) node[pos=0.5,auto=right,inner sep=2pt,text=td1] {$e^-$};
  \end{feynhand}
  \begin{feynhand}
    \vertex[dot,color=td1,line width=0.8pt,minimum size=2mm,inner sep=0pt] (tdv4) at (50,150) {};
  \end{feynhand}
  \begin{feynhand}
    \vertex[dot,color=td1,line width=0.8pt,minimum size=2mm,inner sep=0pt] (tdv5) at (110,150) {};
  \end{feynhand}
  \node[text=td1,anchor=center,inner sep=0pt,outer sep=0pt,align=center,font=\fontsize{11}{13.2}\selectfont] at (80,107) {tikz-feynhand};
  \begin{scope}[american,voltage dir=RP]
    \ctikzset{bipoles/length=12mm}
    \draw[draw=td1,line width=0.8pt] (40,25) to[V,l^={$\text{電源}$}] (40,75);
  \end{scope}
  \begin{scope}[american,voltage dir=RP]
    \ctikzset{bipoles/length=12mm}
    \draw[draw=td0,line width=0.8pt] (40,75) to[R,l^={$R_1$},i>_={$i$}] (115,75);
  \end{scope}
  \begin{scope}[american,voltage dir=RP]
    \ctikzset{bipoles/length=12mm}
    \draw[draw=td1,line width=0.8pt] (115,75) to[C,l_={$C_1$},v^={$v_C$}] (115,25);
  \end{scope}
  \begin{scope}[american,voltage dir=RP]
    \ctikzset{bipoles/length=12mm}
    \draw[draw=td1,line width=0.8pt] (115,25) to[short] (40,25);
  \end{scope}
  \draw[draw=td1,line width=0.8pt] (40,25) node[ground,rotate=0] {};
  \node[circ,color=td1,line width=0.8pt,minimum size=1.6mm,inner sep=0pt] at (40,25) {};
  \node[circ,color=td1,line width=0.8pt,minimum size=1.6mm,inner sep=0pt] at (115,75) {};
  \node[text=td1,anchor=center,inner sep=0pt,outer sep=0pt,align=center,font=\fontsize{11}{13.2}\selectfont] at (80,5) {circuitikz};
\end{tikzpicture}
\end{document}
